Formale Sprachen referieren im Gegensatz zu natürlichen Sprachen auf Sprachen
in der Technik: wie unterhalte ich mich mit einem Rechner oder Roboter? Eins der
wichtigsten Beispiele sind Programmiersprachen auf unterschiedlichstem Niveau,
angefangen von Assembler über gängige Hochsprachen wie Java oder C++ zu
Anwendungssprachen wie etwa Matlab. Andere wichtige Bereiche, in denen
formale Sprachen eine zentrale Rolle spielen, sind etwa XML, oder aber direkte
Maschinenbefehle etwa in der NC Steuerung. Anders als natürliche Sprache muss
dieses so präzise sein, dass die Maschine keinerlei Interpretationsspielraum hat
und genau das versteht, was der Mensch intendiert. Zusätzliche Anforderung ist,
dass dieses auch möglichst effizient umgesetzt werden kann, d.h. der Rechner
sollte ein Wort möglichst schnell erkennen und interpretieren können.


\section*{0\_orga.pdf}
Vorerst irrelevant.


\section{Historisches}
\untit{1\_theoinf\_motivation.pdf - Skript S. 1--2}
Die Formalisierung von Sprache ist ein nicht nur für die Informatik im Rahmen
technischer Systeme relevantes Gebiet, sondern auch im Rahmen der Sprachphi-
losophie und Linguistik behandelt worden. So ist einer der wichtigsten Forscher
in dem Gebiet Noam Chomsky, ein Linguist, der 1956 die nach ihm benannte
Chomsky-Hierarchie von formalen Sprachen einführte (Chomsky 1956). Diese
unterscheidet unterschiedlich komplexe Formalisierungen, die zu unterschiedlich
komplexen Sprachkonstrukten führen. Diese theoretischen Überlegungen hatten in
der Informatik direkte Bedeutung für die Entwicklung von Programmiersprachen.
Die erste Programmiersprache, für die eine formale Syntax eingeführt wurde, war
um 1960 ALGOL 60. Unter anderem war hier John Backus beteiligt, nach dem
die heute noch gebräuchliche Backus-Naur Form von Grammatiken benannt ist.
Basierend auf dieser Formalisierung war die Entwicklung von Compilern deutlich
einfacher und eleganter möglich. \\
Das Gebiet der Formalen Sprachen behandelt unter anderem auch, was auf
unterschiedlichen Komplexitätsstufen darstellbar ist und welche strukturellen
Eigenschaften gelten. Interessanterweise sind einige wichtige strukturelle Eigen-
schaften in dem Gebiet auch noch offen bzw. teilweise sehr spät gelöst worden. Der
Abschluss so genannter kontextsensitiver Sprachen gegen Komplementbildung
wurde erst 1988 bewiesen im Satz von Immerman/Szelepcsenyi (Immerman 1988).
Letzterer war damals noch Student. Er erhielt 1995 zusammen mit Immerman
hierfür den Gödel-Award, den Nobelpreis der Informatik.


\section{Kapitelinhalt}
\untit{2\_0\_sprachen\_intro.pdf - Skript S. 2}
In diesem Kapitel werden wir eine Hierarchie von formalen Sprachen auf der
Basis von Grammatiken und deren Komplexität kennenlernen. Das dient dazu,
unterschiedlich komplexe formale Sprachen definieren und interessante Eigen-
schaften wie etwa den Abschluss untersuchen zu können.
%Wichtige Begriffe in diesem Kapitel sind:


\section{Wortersetzungssysteme und Notationen}
\untit{2\_1\_Sprachen\_Grundlagen.pdf - Skript S. 2--6}
Wortersetzungssysteme oder auch Semi-Thue-Systeme bilden ein grundlegendes
Konzept, Wörter zu erzeugen oder abzuändern. Ausgehend von einem Startsymbol
kann man eine Menge komplexer Wörter durch entsprechende Ersetzungsregeln
definieren. Wir werden Wortersetzungssysteme als Grundlage in der gesamten
Vorlesung verwenden. Zunächst geben wir einige Definitionen:
% 2_0_sprachen_intro.pdf

\defin{symbol}
\defin{alphabet}
\defin{wort}
\defin{anzahl}
\defin{leereswort}
\defin{allewörter}
\defin{allewörterlänge}

\begin{defi}
    \label{defi:Alphabet} \textbf{Alphabet} $\mathbf{\Sigma}$\\
    Endliche Menge von Symbolen:
    \[\Sigma = \{a_1, \ldots, a_n\}\]
    Dabei ist $\Sigma^*$ die Menge aller Wörter über $\Sigma$ und $\Sigma^m$ die Menge aller Wörter der Länge $m$ über $\Sigma$.
\end{defi}
\begin{defi}
    \label{defi:Wort} \textbf{Wort} $\mathbf{\overline{w}}$ über einem Alphabet $\Sigma$ \\
    Eine endliche Folge von Elementen aus einem Alphabet:
    \[\overline{w} = a_{i_1}\ldots a_{i_m}, \quad a_{i_j} \in \Sigma, \quad m > 0\]
    Dabei ist $m = \overline{w}$ die Länge des Wortes und ein Wort der Länge 0 ist das leere Wort \underline{$\epsilon$}
\end{defi}
\begin{defi}
    \label{defi:Anzahl} \textbf{Anzahl} $\mathbf{(a_i, \overline{w})}$ eines Symbols $a_i$ \\
    Die Anzahl der Vorkommnisse des Symbols $a_i$ im
    Wort $\overline{w}$.
\end{defi}
\begin{defi}
    \label{defi:Konkatenation} \textbf{Konkatenation} $\overline{wv}$ zweier Wörter $\overline{w}, \overline{v}$ \\
    Eine Verkettung zweier Wörter $\overline{w} = a_{i_1}\ldots a_{i_m}$ und $\overline{v} = a_{i_1}\ldots a_{i_l}$ zu
    \[\overline{wv}=a_{i_1}\ldots a_{i_m}a_{i_1}\ldots a_{i_l}\]
\end{defi}
\begin{defi}
    \label{defi:Präfix}

\end{defi}

\defin{wortersetzungssystem}
